\documentclass[11pt]{article}

\usepackage[utf8]{inputenc}
\usepackage[margin=1in]{geometry}
\usepackage{amsmath,amssymb,mathtools,enumitem,fancyhdr}
\usepackage[misc]{ifsym}
\usepackage{hyperref}

\setlength\textwidth{7in}
\setlength\parindent{0pt}
\oddsidemargin=-0.5cm
\textheight=665pt
\pagestyle{fancy}
\setlength{\headheight}{13.6pt}
\renewcommand{\headrulewidth}{0.1pt}
\renewcommand{\footrulewidth}{0.1pt}
\lhead{\scshape Pontificia Universidad Católica de Chile}
\rhead{\scshape IMC UC}
\cfoot{\thepage}

\newcommand{\solution}{\large\textbf{Solución}\normalsize}
\newcommand{\norm}[1]{\left\lVert #1\right\rVert}
\newcommand{\abs}[1]{\left|#1\right|}
\def\mat   #1{\mbox{\boldmath $\mathsf #1$}}
\def\vec   #1{\mbox{\boldmath $#1$}{}}
\DeclareMathOperator{\tr}{tr}
\def\R{\mathbb R}
\def\dx{\mbox{\(\,\mathrm{d}x\)}}
\def\ds{\mbox{\(\,\mathrm{d}s\)}}

\AtBeginDocument{\vspace*{-3.2\baselineskip}}

\begin{document}
    \begin{flushleft}
        \small
        \vspace{1pt}
        \textbf{IMT3130} \hfill
        \textbf{10/04/2026}
    \end{flushleft}
    \begin{center}
        \LARGE\textbf{Pauta Ayudantía 4}\\
        \vspace{3pt}
        \normalsize\textbf{Formulaciones débiles, Poincaré, Lax--Milgram y Galerkin}\\
        \normalsize Ayudante: Bastián Herrera \Letter{\hspace{1pt}: \texttt{biherrera@uc.cl}}
        \rule{\textwidth}{0.05mm}
    \end{center}

\section*{Problemas y soluciones}

\begin{enumerate}[label=\textbf{\arabic*.}, leftmargin=*]

    \item \textbf{Enunciado.} Sea $\Omega$ Lipschitz acotado, $\partial\Omega=\overline{\Gamma_D}\cup\overline{\Gamma_R}$, y considere
    \[
    \begin{aligned}
        -\Delta u+cu&=f &&\text{en }\Omega,\\
        \gamma_0u&=0 &&\text{en }\Gamma_D,\\
        \partial_nu+\beta\gamma_0u&=r &&\text{en }\Gamma_R.
    \end{aligned}
    \]
    Con
    \[
        V=\{v\in H^1(\Omega):\gamma_0v=0\text{ en }\Gamma_D\},
    \]
    derive la formulación débil, estudie continuidad/coercividad y analice el caso de Neumann puro.

    \textbf{Solución.}
    \begin{enumerate}[label=(\alph*)]
        \item Tomamos $v\in V$. Multiplicando por $v$ e integrando por partes,
        \begin{align*}
            \int_\Omega fv\dx
            &=
            \int_\Omega (-\Delta u+cu)v\dx\\
            &=
            \int_\Omega \nabla u\cdot\nabla v\dx
            -\int_{\partial\Omega}\partial_nu\,\gamma_0v\ds
            +\int_\Omega cuv\dx.
        \end{align*}
        Como $\gamma_0v=0$ en $\Gamma_D$, sólo queda la contribución de $\Gamma_R$. Usando
        \[
            \partial_nu=r-\beta\gamma_0u\qquad\text{en }\Gamma_R,
        \]
        se obtiene
        \[
            \int_\Omega \nabla u\cdot\nabla v\dx
            +\int_\Omega cuv\dx
            +\int_{\Gamma_R}\beta\,\gamma_0u\,\gamma_0v\ds
            =
            \int_\Omega fv\dx
            +\langle r,\gamma_0v\rangle_{\Gamma_R}.
        \]
        La formulación débil es: hallar $u\in V$ tal que
        \[
            a(u,v)=\ell(v)\qquad\forall v\in V,
        \]
        donde
        \[
            a(w,v)=
            \int_\Omega \nabla w\cdot\nabla v\dx
            +\int_\Omega cwv\dx
            +\int_{\Gamma_R}\beta\,\gamma_0w\,\gamma_0v\ds,
        \]
        y
        \[
            \ell(v)=
            \int_\Omega fv\dx
            +\langle r,\gamma_0v\rangle_{\Gamma_R}.
        \]
        La condición de Dirichlet es esencial porque se incorpora en el espacio de solución y de test. La condición de Robin es natural porque aparece como término de borde después de integrar por partes.

        \item Supongamos, por ejemplo, $f\in L^2(\Omega)$ y $r\in H^{-1/2}(\Gamma_R)$. La continuidad de $a$ sigue de Cauchy--Schwarz, de $c,\beta\in L^\infty$ y de la continuidad de la traza:
        \begin{align*}
            \abs{a(w,v)}
            &\le
            \norm{\nabla w}_{L^2(\Omega)}\norm{\nabla v}_{L^2(\Omega)}
            +\norm{c}_{L^\infty}\norm{w}_{L^2(\Omega)}\norm{v}_{L^2(\Omega)}\\
            &\quad
            +\norm{\beta}_{L^\infty(\Gamma_R)}
            \norm{\gamma_0w}_{L^2(\Gamma_R)}
            \norm{\gamma_0v}_{L^2(\Gamma_R)}\\
            &\le C\norm{w}_{H^1(\Omega)}\norm{v}_{H^1(\Omega)}.
        \end{align*}
        Además,
        \[
            \abs{\ell(v)}
            \le
            \norm{f}_{L^2(\Omega)}\norm{v}_{L^2(\Omega)}
            +\norm{r}_{H^{-1/2}(\Gamma_R)}
            \norm{\gamma_0v}_{H^{1/2}(\Gamma_R)}
            \le C\norm{v}_{H^1(\Omega)}.
        \]

        Si $|\Gamma_D|>0$, la desigualdad de Poincaré generalizada entrega
        \[
            \norm{v}_{H^1(\Omega)}\le C_P\norm{\nabla v}_{L^2(\Omega)}
            \qquad\forall v\in V.
        \]
        Como $c,\beta\ge0$,
        \[
            a(v,v)
            =
            \norm{\nabla v}_{L^2(\Omega)}^2
            +\int_\Omega cv^2\dx
            +\int_{\Gamma_R}\beta(\gamma_0v)^2\ds
            \ge
            \norm{\nabla v}_{L^2(\Omega)}^2
            \ge
            C\norm{v}_{H^1(\Omega)}^2.
        \]
        Luego $a$ es coerciva y Lax--Milgram da existencia y unicidad.

        Si $\Gamma_D=\varnothing$, la coercividad también se recupera si $c(x)\ge c_0>0$ casi en todas partes, pues
        \[
            a(v,v)\ge \norm{\nabla v}_{L^2(\Omega)}^2+c_0\norm{v}_{L^2(\Omega)}^2.
        \]
        También puede recuperarse si el término Robin controla constantes, por ejemplo si $\beta\ge\beta_0>0$ en una porción de borde de medida positiva. Lo que falla en ausencia de estos mecanismos es el control de las constantes.

        \item En Neumann puro,
        \[
            a(w,v)=\int_\Omega \nabla w\cdot\nabla v\dx,
            \qquad
            \ell(v)=\int_\Omega fv\dx+\langle r,\gamma_0v\rangle_{\partial\Omega}.
        \]
        Toda función constante pertenece al núcleo de $a$, porque su gradiente es cero. Si existe solución, al tomar $v=1$ se obtiene necesariamente
        \[
            0=a(u,1)=\ell(1)=\int_\Omega f\dx+\langle r,1\rangle_{\partial\Omega}.
        \]
        Por tanto, la compatibilidad es
        \[
            \int_\Omega f\dx+\langle r,1\rangle_{\partial\Omega}=0.
        \]
        La unicidad sólo puede esperarse módulo constantes. Para fijar una solución única se trabaja en
        \[
            H^1_\ast(\Omega)=\left\{v\in H^1(\Omega):\int_\Omega v\dx=0\right\}.
        \]
        Allí, Poincaré--Wirtinger da
        \[
            \norm{v}_{H^1(\Omega)}\le C\norm{\nabla v}_{L^2(\Omega)},
        \]
        y entonces $a$ es coerciva sobre $H^1_\ast(\Omega)$. Si $c(x)\ge c_0>0$, las constantes ya no están en el núcleo, no se necesita compatibilidad de Neumann puro y el problema es coercivo en todo $H^1(\Omega)$.
    \end{enumerate}

    \item \textbf{Enunciado.} Sea $V$ Hilbert, $a:V\times V\to\R$ simétrica, continua y coerciva, y sea $V_h\subset V$. Si
    \[
        a(u,v)=\ell(v)\quad\forall v\in V,
        \qquad
        a(u_h,v_h)=\ell(v_h)\quad\forall v_h\in V_h,
    \]
    pruebe la ortogonalidad de Galerkin, la equivalencia de la norma de energía, la propiedad de mejor aproximación y escriba el sistema lineal asociado.

    \textbf{Solución.}
    \begin{enumerate}[label=(\alph*)]
        \item Como $u$ y $u_h$ satisfacen las ecuaciones continua y discreta, para todo $v_h\in V_h$,
        \[
            a(u,v_h)=\ell(v_h),
            \qquad
            a(u_h,v_h)=\ell(v_h).
        \]
        Restando,
        \[
            a(u-u_h,v_h)=0\qquad\forall v_h\in V_h.
        \]

        \item La coercividad implica
        \[
            a(v,v)\ge \alpha\norm{v}_V^2,
        \]
        y la continuidad implica
        \[
            a(v,v)\le M\norm{v}_V^2.
        \]
        Por tanto,
        \[
            \sqrt{\alpha}\norm{v}_V\le \norm{v}_a\le \sqrt{M}\norm{v}_V,
        \]
        y $\norm{\cdot}_a$ es una norma equivalente a la norma de $V$.

        \item Sea $v_h\in V_h$ arbitrario. Entonces $v_h-u_h\in V_h$ y por ortogonalidad,
        \[
            a(u-u_h,v_h-u_h)=0.
        \]
        Por tanto,
        \begin{align*}
            \norm{u-v_h}_a^2
            &=
            \norm{(u-u_h)+(u_h-v_h)}_a^2\\
            &=
            \norm{u-u_h}_a^2+\norm{u_h-v_h}_a^2
            +2a(u-u_h,u_h-v_h)\\
            &=
            \norm{u-u_h}_a^2+\norm{u_h-v_h}_a^2
            \ge
            \norm{u-u_h}_a^2.
        \end{align*}
        Como esto vale para todo $v_h\in V_h$ y la igualdad se alcanza con $v_h=u_h$,
        \[
            \norm{u-u_h}_a=\inf_{v_h\in V_h}\norm{u-v_h}_a.
        \]

        \item Si $u_h=\sum_{j=1}^N U_j\varphi_j$, basta testear con cada $\varphi_i$. Se obtiene
        \[
            \sum_{j=1}^N a(\varphi_j,\varphi_i)U_j=\ell(\varphi_i),
            \qquad i=1,\dots,N.
        \]
        En forma matricial,
        \[
            \mat A\vec U=\vec F,
            \qquad
            A_{ij}=a(\varphi_j,\varphi_i),
            \qquad
            F_i=\ell(\varphi_i).
        \]
    \end{enumerate}

\end{enumerate}

\end{document}
